\subsection{INITIAL TOPOLOGIES}

PROPOSITION 4. Let X be a set, let \((\mathbf{Y}_i)_{i \in I}\) be a family of topological spaces, and for each \(i \in I\) let \(f_i\) be a mapping of X into \(\mathbf{Y}_i\) . Let \(\mathfrak{S}\) be the set of subsets of X of the form \(\overline{f}_i^1 (\mathbf{U}_i)\) ( \(i \in I\) , \(\mathbf{U}_i\) open in \(\mathbf{Y}_i\) ), and let \(\mathfrak{B}\) be the set of finite intersections of sets of \(\mathfrak{S}\) . Then \(\mathfrak{B}\) is a base of a topology \(\mathfrak{C}\) on X which is the initial topological structure on X for the family \((f_i)\) (Set Theory, Chapter IV, § 2, no. 3) and in particular is the coarsest topology on X for which the mappings \(f_i\) are continuous. More precisely, if \(g\) is a mapping of a topological space Z into X, then \(g\) is continuous at a point \(z \in Z\) (X carrying the topology \(\mathfrak{C}\) ) if and only if each of the functions \(f_i \circ g\) is continuous at \(z\) .

Let \(\mathfrak{D}\) be the set of all unions of sets belonging to \(\mathfrak{B}\) ; clearly \(\mathfrak{D}\) satisfies axiom \((\mathrm{O}_{\mathrm{I}})\) since formation of unions is associative; and \(\mathfrak{D}\) satisfies axiom \((\mathrm{O}_{\mathrm{II}})\) by reason of the definition of \(\mathfrak{B}\) and the fact that finite intersection is distributive over arbitrary union {[}Set Theory, R § 4, formula (37){]}. The set \(\mathfrak{D}\) is therefore the set of open subsets of X for a topology \(\mathfrak{G}\) of which \(\mathfrak{B}\) is a base. We shall prove the last assertion of the proposition, which implies the others by reason of the general properties of initial structures (Set Theory, Chapter IV, § 2, no. 3, criterion CST 9). In the first place, the definition of \(\mathfrak{S}\) shows that the \(f_{i}\) are continuous on X (no. 1, Theorem 1); hence, if \(g\) is continuous at \(z\) , so are the mappings \(f_{i} \circ g\) (no. 1, Proposition 2). Conversely, suppose that all the mappings \(f_{i} \circ g\) are continuous at \(z\) , and let V be a neighbourhood of \(g(z)\) in X; by definition, there is a finite subset J of I, and for each \(i \in J\) an open subset \(\mathrm{U}_{i}\) of \(\mathrm{Y}_{i}\) such that V contains the set \(\bigcap_{i \in J} f_{i}^{1}(\mathrm{U}_{i})\) and \(g(z)\) belongs to this set. It follows that

\[
\overline {{g}} ^ {1} (\mathrm{V}) \supset \bigcap_ {i \in J} \overline {{g}} ^ {1} (\overline {{f}} _ {i} ^ {1} (\mathrm{U} _ {i})),
\]

and the hypothesis implies that each of the sets \(\overline{g}^{1}(\overline{f}_{i}^{1}(\mathbf{U}_{i}))\) is a neighbourhood of \(z\) in \(\mathbf{Z}\) ; hence \(\overline{g}^{1}(\mathbf{V})\) is also a neighbourhood of \(z\) in \(\mathbf{Z}\) . This completes the proof.

Let \(\mathfrak{B}_i\) be a base of the topology of \(\mathbf{Y}_i(t\in I)\) ; let \(\mathfrak{S}'\) denote the set of subsets of \(X\) of the form \(\overline{f}_i^1 (U_i)\) for \(t\in I\) and \(U_{t}\in \mathfrak{B}_{t}\) for each \(t\in I\) ; if \(\mathfrak{B}'\) is the set of finite intersections of sets of \(\mathfrak{S}'\) , it is evident that \(\mathfrak{B}'\) is a base of the topology \(\mathfrak{C}\) .

The general properties of initial structures (Set Theory, Chapter IV, § 2, no. 3, criterion CST 10) imply in particular the following transitivity property (the direct proof of which is quite straightforward):

PROPOSITION 5. Let X be a set, \((Z_{t})_{t\in I}\) a family of topological spaces, \((J_{\lambda})_{\lambda\in L}\) a partition of I and \((Y_{\lambda})_{\lambda\in L}\) a family of sets indexed by L. Also for each \(\lambda\in L\) let \(h_{\lambda}\) be a mapping of X into \(Y_{\lambda}\) ; for each \(\lambda\in L\) and each \(t\in J_{\lambda}\) let \(g_{t\lambda}\) be a mapping of \(Y_{\lambda}\) into \(Z_{t}\) , and put \(f_{t}=g_{t\lambda}\circ h_{\lambda}\) . If each \(Y_{\lambda}\) carries the coarsest topology for which the mappings \(g_{t\lambda}(t\in J_{\lambda})\) are continuous, then the coarsest topology on X for which the \(f_{t}\) are continuous is the same as the coarsest topology for which the \(h_{\lambda}\) are continuous.

Examples. 1) Inverse image of a topology. Let X be a set, Y a topological space, f a mapping of X into Y; the coarsest topology C on X for which f is continuous is called the inverse image under f of the topology of Y. It follows from Proposition 4 and the formulae for the inverse image of a union and an intersection {[}Set Theory, R, § 4, formulae (34) and (46){]} that the open (resp. closed) sets in the topology C are the inverse images under f of the open (resp. closed) sets of Y; consequently, for each \(x \in X\) , the sets \(\overline{f}^{1}(W)\) , where W runs through a fundamental system of neighbourhoods of \(f(x)\) in Y, form a fundamental system of neighbourhoods of x in the topology C. In § 3 we shall study, under the name of induced topology, the particular case in which X is a subset of Y and f is the canonical injection \(X \to Y\) ; X, with the induced topology, is then called a subspace of Y.

For a mapping \(f\) of a topological space \(\mathbf{X}\) into a topological space \(\mathbf{X}'\) to be continuous it is necessary and sufficient that the topology of \(\mathbf{X}\) is finer than the inverse image under \(f\) of the topology of \(\mathbf{X}'\) .

\begin{enumerate}
\def\labelenumi{\arabic{enumi})}
\setcounter{enumi}{1}
\tightlist
\item
  Least upper bound of a set of topologies. Every family \((\mathcal{C}_t)_{i\in I}\) of topologies on a set \(X\) has a least upper bound \(\mathcal{C}\) in the ordered set of all topologies on \(X\) , i.e.~there exists a topology on \(X\) which is coarsest among all the topologies on \(X\) which are finer than each of the \(\mathcal{C}_t\) . To see this we may apply Proposition 4, taking \(Y_t\) to be the set \(X\) with the topology \(\mathcal{C}_t\) , and \(f_t\) to be the identity map \(X\to Y_t\) ; \(\mathcal{C}\) is the coarsest topology for which all the mappings \(f_t\) are continuous.
\end{enumerate}

Let \(\mathfrak{S}\) be an arbitrary set of subsets of a set X; amongst the topologies \(\mathfrak{C}\) on X for which the sets of \(\mathfrak{S}\) are open, there is a topology \(\mathfrak{C}_0\) which is coarser than all the others and which is called the topology generated by \(\mathfrak{S}\) . For each set \(U \in \mathfrak{S}\) let \(\mathfrak{C}_U\) be the topology whose open sets are \(\varnothing\) , U and X {[}it is clear that this set of subsets of X satisfies \((O_I)\) and \((O_{II})]\) ; then \(\mathfrak{C}_0\) is just the least upper bound of the topologies \(\mathfrak{C}_U\) . By Proposition 4, if \(\mathfrak{B}\) is the set of finite intersections of sets belonging to \(\mathfrak{S}\) , then \(\mathfrak{B}\) is a base of the topology \(\mathfrak{C}_0\) . We say that \(\mathfrak{S}\) is a sub-base of \(\mathfrak{C}_0\) .

\begin{enumerate}
\def\labelenumi{\arabic{enumi})}
\setcounter{enumi}{2}
\tightlist
\item
  Product topology. Let \((\mathbf{X}_t)_{i \in I}\) be a family of topological spaces. The coarsest topology on the product set \(\mathbf{X} = \prod_{i \in I} \mathbf{X}_i\) for which the projec-
\end{enumerate}

tions \(pr_{i}: X \rightarrow X_{i}\) are continuous mappings is called the product of the topologies of the \(X_{i}\) ; we shall study it in more detail in § 4.

